\documentclass[a4paper, serif, 10pt]{moderncv}

%% moderncv
% color and style
\moderncvcolor{black}
\moderncvstyle{banking}

%% page layout/margins
\usepackage[top=2.5cm, bottom=2.5cm, left=1.5cm, right=1.5cm]{geometry}

\nopagenumbers{}

%% Babel config for Spanish language
% \usepackage[spanish]{babel} has some conflicting issues with moderncv
% so we have to use enable the following options to make it work
%
% ref:
% - https://tex.stackexchange.com/a/140161/36007
\usepackage[spanish,es-lcroman]{babel}

%% fontspec
\usepackage{fontspec}

\IfFontExistsTF{Linux Libertine}{
  \setmainfont[Ligatures={TeX, Common}, Numbers=OldStyle]{Linux Libertine}
}{}
\IfFontExistsTF{Linux Libertine O}{
  \setmainfont[Ligatures={TeX, Common}, Numbers=OldStyle]{Linux Libertine O}
}{}

%% CTeX
% CJK support, used to show CJK characters in the resume
%
% - heading=false: disable ctex heading
% - punct=kaiming: use kaiming punctuations styles for CJK
% - scheme=plain: use plain scheme, do not override \normalsize font size
% - space=auto: space settings for CJK characters
%
% CJK fonts are resolved per locale language, e.g. Noto Serif/Sans CJK JP
% for Japanese resumes, with fallback to the next region if not installed
%
% ref:
% - http://ctan.mirrorcatalogs.com/language/chinese/ctex/ctex.pdf
\usepackage[UTF8, heading=false, punct=kaiming, scheme=plain, space=auto]{ctex}

\IfFontExistsTF{Noto Serif CJK SC}{
  \setCJKmainfont{Noto Serif CJK SC}
}{}
\IfFontExistsTF{Noto Sans CJK SC}{
  \setCJKsansfont{Noto Sans CJK SC}
}{}

%% line spacing
\usepackage{setspace}
\setstretch{1.125}

%% URL styling - use normal text instead of monospace
\urlstyle{same}

%% Patch \httplink and \httpslink to handle full URLs
% The original moderncv macros always prepend http:// or https:// to URLs.
% This causes issues when the URL already has a protocol (e.g., https://example.com)
% resulting in malformed URLs like https://https://example.com.
% This patch checks if the URL already contains "://" and uses it directly if so.
% Note: We use \str_set:Nx (x-expansion) to expand macros like \@homepage first.
\ExplSyntaxOn
\str_new:N \g_yamlresume_url_str

\RenewDocumentCommand{\httpslink}{O{}m}{
  \str_set:Nx \g_yamlresume_url_str {#2}
  \str_if_in:NnTF \g_yamlresume_url_str {://}
    { \url{#2} }
    { \url{https://#2} }
}

\RenewDocumentCommand{\httplink}{O{}m}{
  \str_set:Nx \g_yamlresume_url_str {#2}
  \str_if_in:NnTF \g_yamlresume_url_str {://}
    { \url{#2} }
    { \url{http://#2} }
}
\ExplSyntaxOff

\name{Lucía Fernández Gómez}{}
\title{Ingeniera QA Senior | Automatización de Pruebas y Calidad de Software}
\phone[mobile]{+34 612 345 678}
\email{lucia.fernandez.qa@example.com}
\homepage{https://luciafernandezqa.dev}

\address{Calle de Alcalá 120, 4º B, Madrid, Comunidad de Madrid, España, 28009}{}{}

\extrainfo{{\small \faLinkedin}\ \href{https://www.linkedin.com/in/lucia-fernandez-qa}{@lucia-fernandez-qa} {} {} {} • {} {} {} 
{\small \faGithub}\ \href{https://github.com/luciafdzqa}{@luciafdzqa}}

\begin{document}

\maketitle

\section{Información básica}

\cvline{}{\begin{itemize}
\item Ingeniera de calidad de software con más de 6 años de experiencia diseñando estrategias de prueba para aplicaciones web y APIs en entornos ágiles.
\item Especializada en automatización E2E y de API con Playwright, Cypress, TypeScript y Python, integrando suites en pipelines CI/CD.
\item Experiencia liderando equipos QA, definiendo métricas de calidad y promoviendo la cultura de prevención de defectos.
\item Colaboradora cercana con producto, desarrollo y operaciones para asegurar entregas rápidas, estables y centradas en el usuario.
\end{itemize}}
\section{Educación}

\cventry{sept 2015–jul 2019}
        {Licenciatura, Ingeniería Informática, Puntuación: 8.7/10}
        {Universidad Politécnica de Madrid}
        {\url{https://www.upm.es}}
        {}
        {\begin{itemize}
\item Formación sólida en desarrollo de software, bases de datos y sistemas distribuidos.
\item Proyecto final orientado a la validación automatizada de servicios REST.
\item Participación en equipos multidisciplinares y prácticas en laboratorio de calidad.
\end{itemize}
\textbf{Cursos}: Ingeniería del Software, Bases de Datos, Sistemas Operativos, Redes de Computadores, Arquitectura de Software, Seguridad Informática, Metodologías Ágiles, Calidad y Pruebas de Software}
\section{Experiencia laboral}

\cventry{mar 2022 hasta la fecha}
        {Senior QA Automation Engineer}
        {NexoPay Tecnología}
        {\url{https://www.nexopay.es}}
        {}
        {\begin{itemize}
\item Diseñé la estrategia de pruebas para una plataforma de pagos con más de 2 millones de transacciones mensuales.
\item Implementé suites de automatización E2E con Playwright y TypeScript, reduciendo el tiempo de regresión un 65 \%.
\item Desarrollé pruebas de API y contratos con Postman, Newman y OpenAPI, integradas en GitHub Actions.
\item Lideré un equipo de tres ingenieros QA, fomentando buenas prácticas de revisión de código y documentación.
\item Definí indicadores de calidad y dashboards en Grafana para seguimiento de defectos, cobertura y estabilidad.
\end{itemize}
\textbf{Palabras clave}: Playwright, TypeScript, Pruebas de API, GitHub Actions, Liderazgo QA, Pagos}

\cventry{jun 2019–feb 2022}
        {QA Engineer}
        {Grupo Aire Tecnología}
        {\url{https://www.grupoairetecnologia.com}}
        {}
        {\begin{itemize}
\item Automaticé pruebas funcionales y de regresión para aplicaciones web con Selenium WebDriver y Python.
\item Gestioné casos de prueba, defectos y trazas en Jira y TestRail dentro de equipos Scrum.
\item Ejecuté pruebas de integración, compatibilidad y accesibilidad en distintos navegadores y dispositivos.
\item Participé en revisiones de requisitos y criterios de aceptación para prevenir defectos desde el diseño.
\item Colaboré con desarrollo en la resolución de incidencias críticas y en la mejora continua del proceso de entrega.
\end{itemize}
\textbf{Palabras clave}: Selenium, Python, Jira, TestRail, Scrum, Regresión}

\cventry{sept 2018–may 2019}
        {QA Analyst Junior}
        {Soluciones Móviles Ibéricas}
        {\url{https://www.solucionesmovilesibericas.com}}
        {}
        {\begin{itemize}
\item Ejecuté pruebas manuales funcionales y exploratorias sobre aplicaciones móviles iOS y Android.
\item Documenté evidencias, reproducibilidad y severidad de defectos en herramientas de seguimiento.
\item Apoyé la creación de checklists y matrices de dispositivos para pruebas de compatibilidad.
\item Aprendí fundamentos de ciclo de vida del defecto y metodologías ágiles en un entorno de producto.
\end{itemize}
\textbf{Palabras clave}: Pruebas manuales, Mobile QA, Android, iOS, Checklists, Defectos}
\section{Idiomas}

\cvline{Español}{Competencia nativa o bilingüe \hfill \textbf{Palabras clave}: Nativo}
\cvline{Inglés}{Competencia profesional plena \hfill \textbf{Palabras clave}: TOEIC 940, Comunicación técnica}
\cvline{Portugués}{Competencia limitada de trabajo \hfill \textbf{Palabras clave}: Conversación}
\section{Competencias}

\cvline{Automatización de Pruebas}{Experto \hfill \textbf{Palabras clave}: Playwright, Cypress, Selenium, TypeScript, Python}
\cvline{Pruebas de API}{Avanzado \hfill \textbf{Palabras clave}: Postman, Newman, REST Assured, OpenAPI, GraphQL}
\cvline{CI/CD y DevOps}{Avanzado \hfill \textbf{Palabras clave}: GitHub Actions, Jenkins, Docker, Git, Azure DevOps}
\cvline{Gestión de Calidad}{Avanzado \hfill \textbf{Palabras clave}: Jira, TestRail, Xray, Métricas QA, Gestión de defectos}
\cvline{Pruebas de Rendimiento}{Intermedio \hfill \textbf{Palabras clave}: JMeter, k6, Lighthouse, Gatling}
\section{Premios}

\cventry{nov 2023}
        {NexoPay Tecnología}
        {Premio a la Excelencia en Calidad}
        {}
        {}
        {Reconocimiento por liderar la modernización de la estrategia de pruebas y reducir un 65 \% el tiempo de regresión.}

\cventry{jun 2021}
        {Grupo Aire Tecnología}
        {Reconocimiento a la Innovación en Automatización}
        {}
        {}
        {Otorgado por el diseño de un framework interno de pruebas de API reutilizable por tres equipos de producto.}
\section{Certificados}

\cventry{abr 2021}
        {ISTQB}
        {ISTQB Certified Tester Advanced Level - Test Automation Engineer}
        {\url{https://www.istqb.org/certifications/advanced-level-test-automation-engineer}}
        {}
        {}

\cventry{feb 2023}
        {Postman}
        {Postman API Fundamentals Student Expert}
        {\url{https://www.postman.com/postman/workspace/postman-student-expert-program}}
        {}
        {}

\cventry{sept 2022}
        {Amazon Web Services}
        {AWS Certified Cloud Practitioner}
        {\url{https://aws.amazon.com/certification/certified-cloud-practitioner/}}
        {}
        {}
\section{Publicaciones}

\cventry{may 2023}
        {Pruebas de contrato en arquitecturas de microservicios}
        {Revista Tecnología y Calidad}
        {\url{https://www.revistatecnologiaycalidad.com/articulos/pruebas-contrato-microservicios}}
        {}
        {Artículo técnico sobre estrategias de pruebas de contrato con Pact y OpenAPI para equipos que migran a microservicios.}
\section{Proyectos}

\cventry{ene 2023–ago 2023}
        {Framework interno para pruebas E2E y de API reutilizable en varios equipos.}
        {Framework de Automatización QA}
        {\url{https://github.com/luciafdzqa/qa-automation-framework}}
        {}
        {\begin{itemize}
\item Diseñé una arquitectura modular con Playwright, TypeScript y Page Objects.
\item Integré reportes HTML, capturas y trazas en GitHub Actions.
\item Documenté guías de contribución y estándares de codificación para QA.
\end{itemize}
\textbf{Palabras clave}: Playwright, TypeScript, Page Objects, CI/CD, Framework}

\cventry{mar 2022–nov 2022}
        {Iniciativa para incorporar validaciones de accesibilidad en el ciclo de pruebas.}
        {Programa de Pruebas de Accesibilidad}
        {\url{https://github.com/luciafdzqa/accesibilidad-qa}}
        {}
        {\begin{itemize}
\item Definí criterios WCAG 2.1 AA y checklists para equipos de producto.
\item Automaticé auditorías con axe-core y Lighthouse en pipelines de integración.
\item Formé a desarrolladores y diseñadores en buenas prácticas de accesibilidad.
\end{itemize}
\textbf{Palabras clave}: Accesibilidad, WCAG, axe-core, Lighthouse, Calidad}
\section{Intereses}

\cvline{Tecnología}{Automatización, Open Source, Ciberseguridad}
\cvline{Deporte}{Escalada, Running, Ciclismo}
\cvline{Música}{Guitarra, Conciertos, Producción musical}
\section{Voluntariado}

\cventry{sept 2021 hasta la fecha}
        {Mentora y Organizadora}
        {Comunidad QA Madrid}
        {\url{https://www.comunidadqamadrid.org}}
        {}
        {\begin{itemize}
\item Organizo encuentros mensuales sobre automatización, calidad y liderazgo técnico.
\item Mentoro a personas que inician su carrera en QA y reviso sus primeros proyectos prácticos.
\item Coordino talleres de pruebas de API y accesibilidad con ponentes de la industria.
\end{itemize}}
    
\end{document}